\chapter{Peluang dan Peubah Acak}\label{ch:g11:prob}

Peluang memodelkan percobaan yang hasilnya tidak pasti: lemparan dadu, kartu
yang diambil, permainan yang dimainkan. Bab ini menyiapkan kosakatanya pada
himpunan berhingga, lalu memperkenalkan tokoh utama seluruh teori peluang:
\emph{\omterm{def:g11:prob:rv}{peubah acak}} dan \emph{\omterm{def:g11:prob:expectation}{nilai harapannya}}. Gagasan ini ditinjau
ulang dan diperdalam pada \cref{ch:g12:randvar}, dan peluang bersyarat
menyusul pada \cref{ch:g12:condprob}.

\section{Peluang pada himpunan berhingga}

\begin{definition}[Model peluang]\label{def:g11:prob:model}
Percobaan dengan berhingga banyak hasil dimodelkan oleh \emph{ruang
sampelnya}\index{ruang sampel} $\Omega = \{\omega_1, \dots, \omega_n\}$
(himpunan hasilnya) beserta sebuah \emph{peluang} $\P$ yang memberikan kepada
setiap hasil $\omega_i$ sebuah bilangan $p_i \geq 0$, dengan
$p_1 + \dots + p_n = 1$. Sebuah \emph{kejadian}\index{kejadian} adalah
himpunan bagian $A \subseteq \Omega$, dan
\[
\P(A) = \sum_{\omega_i \in A} p_i .
\]
Ketika semua hasilnya sama-sama mungkin ($p_i = \frac1n$: \emph{model
seragam}\index{model seragam}),
\[
\P(A) = \frac{\text{banyaknya hasil dalam } A}
{\text{banyaknya hasil dalam } \Omega}.
\]
\end{definition}

\begin{example}\label{ex:g11:prob:twodice}
Lemparlah dua dadu yang dapat dibedakan: $\Omega$ adalah himpunan $36$ pasangan
terurut $(i, j)$ dengan $1 \leq i, j \leq 6$, semuanya sama-sama mungkin.
\omterm{def:g11:prob:model}{Kejadian} $S$ = ``jumlahnya $7$'' memuat $6$ pasangan
$(1,6)$, $(2,5)$, $(3,4)$, $(4,3)$, $(5,2)$, $(6,1)$, sehingga
$\P(S) = \frac{6}{36} = \frac16$. \omterm{def:g11:prob:model}{Kejadian} ``jumlahnya $12$'' hanya memuat
$(6,6)$: peluangnya $\frac{1}{36}$. Memilih \omterm{def:g11:prob:model}{ruang sampel} yang \emph{tepat} ---
pasangan terurut, bukan takterurut --- itulah yang membuat \omterm{def:g11:prob:model}{model
seragamnya} dapat diterapkan.
\end{example}

\begin{proposition}[Aturan peluang]\label{prop:g11:prob:rules}
Untuk \omterm{def:g11:prob:model}{kejadian} $A, B \subseteq \Omega$:
\[
\P(\emptyset) = 0, \qquad \P(\Omega) = 1, \qquad
\P(\bar A) = 1 - \P(A),
\]
\[
\P(A \cup B) = \P(A) + \P(B) - \P(A \cap B),
\]
dengan $\bar A$ adalah \emph{\omterm{def:g10:proba:operations}{komplemen}} $A$ (semua hasil yang tidak dalam $A$).
Khususnya, $\P(A \cup B) = \P(A) + \P(B)$ ketika $A$ dan $B$
\emph{\omterm{def:g10:proba:operations}{saling lepas}} ($A \cap B = \emptyset$).
\end{proposition}

\begin{proof}
\omterm{def:g11:prob:model}{Kejadian} kosong tidak memuat hasil apa pun (jumlah kosong), dan $\Omega$ memuat
semuanya (jumlah totalnya $1$). Setiap hasil berada tepat pada salah satu dari
$A$ dan $\bar A$, sehingga $\P(A) + \P(\bar A) = 1$. Untuk gabungannya:
menjumlahkan $p_i$ atas $A$ lalu atas $B$ membilang hasil dalam $A \cap B$ dua
kali; mengurangkan $\P(A \cap B)$ membetulkan penghitungan gandanya.
\end{proof}

\begin{example}
Di sebuah kelas, $60\%$ belajar musik ($M$), $45\%$ berolahraga ($S$), dan
$25\%$ melakukan keduanya. Bagian yang melakukan paling sedikit satu di
antaranya adalah $\P(M \cup S) = 0.60 + 0.45 - 0.25 = 0.80$, sehingga $20\%$
tidak melakukan keduanya (\omterm{def:g10:proba:operations}{komplemennya}).
\end{example}

\section{Peubah acak}

\begin{definition}[Peubah acak, distribusi]\label{def:g11:prob:rv}
Sebuah \emph{peubah acak}\index{peubah acak} pada $\Omega$ adalah \omterm{def:g11:func:function}{fungsi}
$X$ yang memberikan sebuah bilangan kepada setiap hasil.
Yang disebut \emph{distribusi}\index{distribusi} peubah itu adalah tabel nilai yang mungkin
$x_1, \dots, x_k$ beserta peluangnya
\[
\P(X = x_i) = \P\bigl(\{\omega : X(\omega) = x_i\}\bigr),
\qquad
\sum_{i=1}^k \P(X = x_i) = 1 .
\]
\end{definition}

\begin{example}\label{ex:g11:prob:game}
Sebuah permainan: bayar $2$ euro, lempar satu dadu setimbang, terimalah nilai
mata yang muncul bila paling sedikit $5$, dan tidak menerima apa-apa selain
itu. Misalkan $G$ \emph{keuntungannya} (yang diterima dikurangi yang dibayar).
Hasil $1, 2, 3, 4$ memberi $G = -2$; hasil $5$
memberi $G = 3$; hasil $6$ memberi $G = 4$. \omterm{def:g11:prob:rv}{Distribusi} $G$:
\begin{center}
\begin{tabular}{c|ccc}
$g$ & $-2$ & $3$ & $4$\\
\hline\rule{0pt}{11pt}%
$\P(G = g)$ & $\dfrac46$ & $\dfrac16$ & $\dfrac16$
\end{tabular}
\end{center}
\end{example}

\begin{method}[Menyusun tabel distribusi]\label{met:g11:prob:table}
Daftarkan nilai $X$ yang mungkin; untuk setiap nilai, kumpulkan hasil yang
menghasilkannya lalu jumlahkan peluangnya; periksalah bahwa baris
peluangnya berjumlah $1$.
\end{method}

\section{Nilai harapan}

\begin{definition}[Nilai harapan]\label{def:g11:prob:expectation}
\emph{Nilai harapan}\index{nilai harapan} (atau rata-rata) $X$ adalah rata-rata
nilainya yang dibobot oleh peluangnya:
\[
\E(X) = \sum_{i=1}^{k} \P(X = x_i)\; x_i .
\]
\end{definition}

\begin{remark}
$\E(X)$ adalah rata-rata yang akan diperoleh dari pengulangan percobaan itu
yang sangat banyak dan saling bebas --- versi tepat pernyataan ini adalah
hukum bilangan besar, yang dibuktikan pada \cref{ch:g12:sums}. Perhatikan
keserupaan bentuknya dengan rata-rata sebuah kumpulan data berfrekuensi
(\cref{def:g11:stat:mean}): peluang berperan sebagai \omterm{def:g10:stats:series}{frekuensi
relatif} $\frac{n_i}{n}$.
\end{remark}

\begin{example}\label{ex:g11:prob:gameexp}
Untuk permainan pada \cref{ex:g11:prob:game}:
\[
\E(G) = \frac46 \times (-2) + \frac16 \times 3 + \frac16 \times 4
= \frac{-8 + 3 + 4}{6} = -\frac16 \approx -0.17 .
\]
Secara rata-rata pemainnya rugi kira-kira $17$ sen setiap permainan: jadi
permainan itu merugikan. Sebuah permainan disebut \emph{adil} bila harapan keuntungannya $0$.
\end{example}

\begin{proposition}[Kelinearan]\label{prop:g11:prob:linearity}
Untuk semua \omterm{def:g10:numbers:sets}{bilangan real} $a, b$:
\[
\E(aX + b) = a\,\E(X) + b .
\]
\end{proposition}

\begin{proof}
Peubah $aX + b$ mengambil nilai $a x_i + b$ dengan peluang
$\P(X = x_i)$, sehingga
\[
\E(aX + b) = \sum_i \P(X = x_i)(a x_i + b)
= a \sum_i \P(X = x_i)\,x_i + b \underbrace{\sum_i \P(X =
x_i)}_{=\,1} = a\,\E(X) + b .
\]
\end{proof}

\section{Ragam dan simpangan baku}

\begin{definition}[Ragam]\label{def:g11:prob:variance}
\emph{Ragam}\index{ragam} $X$ mengukur sebaran nilainya
di sekitar \omterm{def:g11:prob:expectation}{nilai harapan} $m = \E(X)$:
\[
\V(X) = \E\bigl((X - m)^2\bigr)
= \sum_{i=1}^{k} \P(X = x_i)\,(x_i - m)^2,
\qquad
\sigma(X) = \sqrt{\V(X)} .
\]
\end{definition}

\begin{proposition}[Rumus jalan pintas]\label{prop:g11:prob:konig}
\[
\V(X) = \E(X^2) - \E(X)^2 .
\]
\end{proposition}

\begin{proof}
Jabarkan $(x_i - m)^2 = x_i^2 - 2m\,x_i + m^2$ di dalam jumlahnya:
\[
\V(X) = \sum_i \P(X = x_i)\,x_i^2
- 2m \sum_i \P(X = x_i)\,x_i + m^2\sum_i \P(X = x_i)
= \E(X^2) - 2m^2 + m^2 ,
\]
yang sama dengan $\E(X^2) - m^2$. Inilah versi \omterm{def:g11:prob:rv}{peubah acak} dari
identitas yang dibuktikan untuk kumpulan data pada \cref{prop:g11:stat:shortcut}.
\end{proof}

\begin{proposition}[Transformasi afin]\label{prop:g11:prob:affine}
$\V(aX + b) = a^2\,\V(X)$ dan $\sigma(aX + b) = \abs a\,\sigma(X)$.
\end{proposition}

\begin{proof}
Menurut kelinearannya, $aX + b$ menyimpang dari \omterm{def:g11:prob:expectation}{nilai harapannya} $a m + b$
sebesar $(aX + b) - (am + b) = a(X - m)$: jadi setiap simpangannya terskalakan
dengan $a$, sehingga setiap kuadrat simpangannya terskalakan dengan $a^2$, dan
begitu pula rata-rata berbobotnya. Pergeseran $b$ lenyap --- menggeser sebuah
peubah tidak mengubah sebarannya.
\end{proof}

\begin{example}\label{ex:g11:prob:variance}
Untuk permainan pada \cref{ex:g11:prob:game}:
$\E(G^2) = \frac46 \times 4 + \frac16 \times 9 + \frac16 \times 16
= \frac{16 + 9 + 16}{6} = \frac{41}{6}$, sehingga
\[
\V(G) = \frac{41}{6} - \left(-\frac16\right)^2
= \frac{41}{6} - \frac{1}{36} = \frac{245}{36} \approx 6.8,
\qquad
\sigma(G) \approx 2.6 .
\]
Kerugian $0.17$ per permainan secara rata-rata itu datang bersama ayunan
sebesar $2.6$ euro --- jadi \omterm{def:g11:prob:expectation}{nilai harapan} tanpa \omterm{def:g11:stat:variance}{simpangan baku} hanya
menceritakan separuh kisahnya.
\end{example}

\begin{omfigure}
\begin{tikzpicture}
\begin{axis}[
  width=8.6cm, height=5.2cm,
  ybar, bar width=14pt,
  xlabel={$g$}, ylabel={$\P(G = g)$},
  xmin=-3.5, xmax=5.5, ymin=0, ymax=0.8,
  xtick={-2,3,4}, ytick={1/6, 4/6},
  yticklabels={$\tfrac16$,$\tfrac46$},
  tick label style={font=\small},
  label style={font=\small}]
  \addplot[fill=omDef!40, draw=omDef] coordinates {
    (-2,0.6667) (3,0.1667) (4,0.1667)};
  \draw[omThm, dashed, thick] (axis cs:-0.1667,0) --
    (axis cs:-0.1667,0.75);
  \node[omThm, font=\small, anchor=south west] at (axis cs:-0.05,0.62)
    {$\E(G)$};
\end{axis}
\end{tikzpicture}

{\small \omterm{def:g11:prob:rv}{Distribusi} keuntungan $G$ pada
\cref{ex:g11:prob:game}, dan \omterm{def:g11:prob:expectation}{nilai harapannya} $\E(G) = -\frac16$ (putus-putus):
yaitu titik seimbang massa peluangnya.}
\end{omfigure}

\section{Latihan}

\begin{exercise}[$\star$]\label{exo:g11:prob:1}
Sehelai kartu diambil dari satu set $52$ kartu baku. Hitunglah peluang
mengambil: sehelai hati; seorang raja; sehelai hati atau seorang raja; bukan
hati dan bukan pula raja.
\end{exercise}

\begin{exercise}[$\star$]\label{exo:g11:prob:2}
Dua dadu setimbang dilempar dan jumlah matanya $S$ dicatat. Dengan
model $36$ hasil pada \cref{ex:g11:prob:twodice}, hitunglah $\P(S = 6)$,
$\P(S \geq 10)$ dan $\P(S \text{ genap})$.
\end{exercise}

\begin{exercise}[$\star$]\label{exo:g11:prob:3}
Sebuah \omterm{def:g11:prob:rv}{peubah acak} mempunyai \omterm{def:g11:prob:rv}{distribusi}
\begin{center}
\begin{tabular}{c|cccc}
$x$ & $-1$ & $0$ & $2$ & $5$\\
\hline
$\P(X=x)$ & $0.3$ & $0.2$ & $0.4$ & $p$
\end{tabular}
\end{center}
Carilah $p$, lalu hitunglah $\E(X)$.
\end{exercise}

\begin{exercise}[$\star$]\label{exo:g11:prob:4}
$X$ memenuhi $\E(X) = 3$ dan $\V(X) = 4$. Berikan \omterm{def:g11:prob:expectation}{nilai harapan}, \omterm{def:g11:prob:variance}{ragam}
dan \omterm{def:g11:stat:variance}{simpangan baku} $Y = 2X - 5$.
\end{exercise}

\begin{exercise}[$\star\star$]\label{exo:g11:prob:5}
Sebuah roda keberuntungan mempunyai $10$ juring yang sama: $5$ bertuliskan $0$
euro, $3$ bertuliskan $2$ euro, $2$ bertuliskan $5$ euro. Bermain berbiaya $2$
euro. Berikan \omterm{def:g11:prob:rv}{distribusi} keuntungan $G$ (perolehan dikurangi taruhan),
hitunglah $\E(G)$, lalu putuskan apakah permainannya menguntungkan.
\end{exercise}

\begin{exercise}[$\star\star$]\label{exo:g11:prob:6}
Sebuah guci berisi $3$ bola merah dan $2$ bola biru; dua bola diambil
\emph{tanpa} pengembalian. Misalkan $X$ banyaknya bola merah yang terambil.
Susunlah \omterm{def:g11:prob:rv}{distribusi} $X$ (daftarkan $20$ pengambilan terurutnya, atau pakailah
sebuah \omterm{met:g10:proba:tree}{pohon}), lalu hitunglah $\E(X)$.
\end{exercise}

\begin{exercise}[$\star\star$]\label{exo:g11:prob:7}
Untuk roda pada \cref{exo:g11:prob:5}, hitunglah $\V(G)$ dan $\sigma(G)$.
\end{exercise}

\begin{exercise}[$\star\star$]\label{exo:g11:prob:8}
Sebuah perusahaan asuransi menjual polis seharga $80$ euro. Dengan peluang
$0.05$ nasabahnya mengeklaim $1000$ euro, dengan peluang $0.01$ nasabahnya
mengeklaim $5000$ euro, dan selain itu tidak ada klaim. Misalkan $B$ manfaat
yang diperoleh perusahaan dari satu polis. Berikan \omterm{def:g11:prob:rv}{distribusi} $B$, hitunglah
$\E(B)$, lalu tafsirkan.
\end{exercise}

\begin{exercise}[$\star\star$]\label{exo:g11:prob:9}
Sebuah soal pilihan ganda mempunyai $4$ jawaban, satu di antaranya benar.
Jawaban benar bernilai $+1$; untuk mencegah terkaan, jawaban salah bernilai
$x < 0$. Seorang siswa menjawab secara acak. Untuk nilai $x$ yang mana harapan
nilai siswa itu sama dengan nol?
\end{exercise}

\begin{exercise}[$\star\star$]\label{exo:g11:prob:10}
Dua koin setimbang dilempar. Ani mengusulkan: ``jika kedua koinnya sama, saya
membayarmu $1$ euro; jika berbeda, kamu membayarku $1$ euro.'' Hitunglah harapan
keuntungan setiap pemain. Apakah permainannya adil? Pertanyaan yang sama dengan
tiga koin, dengan Ani membayar $2$ euro bila ketiganya sama.
\end{exercise}

\begin{exercise}[$\star\star\star$]\label{exo:g11:prob:11}
Sebuah undian menjual $1000$ karcis seharga $5$ euro masing-masing. Satu karcis
memenangkan $2000$ euro, lima karcis memenangkan $200$ euro, dan dua puluh
karcis memenangkan $25$ euro.
\begin{enumerate}
  \item Berikan \omterm{def:g11:prob:rv}{distribusi} pendapatan total penyelenggara dikurangi
        hadiahnya, dan \omterm{def:g11:prob:rv}{distribusi} keuntungan $G$ seorang pemain.
  \item Hitunglah $\E(G)$ dan total harapan laba penyelenggaranya.
        Periksalah bahwa kedua jawabannya taat asas.
\end{enumerate}
\end{exercise}

\section{Soal: Permainan yang tak selesai}

\begin{problem}[{Soal akhir pekan --- Pascal, Fermat dan pembagian
adil sebuah taruhan yang terhenti: kelahiran nilai harapan, keajaiban
kelinearan, dan apa yang tak terlihat oleh nilai harapan}]
\label{pb:g11:prob:1}
Pada tahun 1654 Chevalier de Méré --- penjudi yang sama, yang taruhan dadunya
membuka soal peluang di jilid sebelumnya ---
mengajukan pertanyaan yang lebih dalam kepada Blaise Pascal: dua
pemain yang sama mahirnya harus menghentikan rangkaian permainannya pada
kedudukan $5$ lawan $3$, sedangkan enam kemenangan mengambil taruhan
$64$ pistole. \emph{Bagaimana taruhan itu harus dibagi secara adil?}
Surat-menyurat antara Pascal dan Fermat yang menjawabnya
mendirikan gagasan yang menjadi inti bab ini: \omterm{def:g11:prob:expectation}{nilai harapan}
(\cref{def:g11:prob:expectation}). Kamu akan menurunkan ulang
jawaban mereka dengan kedua cara, lalu berjumpa dengan sulap kelinearan dan
satu titik buta \omterm{def:g11:prob:expectation}{nilai harapan}.

\medskip
\noindent\textbf{Bagian I --- Kefasihan dengan \omterm{def:g11:prob:expectation}{nilai harapan}.}
\begin{enumerate}
  \item Lemparlah dadu setimbang; keuntunganmu adalah (mata dadu $- 3$) euro.
        Susunlah tabel \omterm{def:g11:prob:rv}{distribusi} keuntungan $G$
        (\cref{met:g11:prob:table}) lalu hitunglah $\E(G)$.
  \item Hitunglah $V(G)$ dengan rumus jalan pintasnya
        (\cref{prop:g11:prob:konig}) dan $\sigma(G)$ (dua
        angka desimal).
  \item Tanpa tabel baru apa pun, berikan $\E(2G - 1)$ dan
        $V(2G - 1)$ (\cref{prop:g11:prob:affine}).
  \item Sebuah lapak menarik $m$ euro untuk melempar dua dadu lalu
        menyerahkan jumlah matanya kepadamu dalam euro. Berapa harga
        $m$ yang \emph{adil}? Lapak itu menarik $8$: berapa harapan kerugianmu
        per permainan?
  \item Sebuah laptop seharga $800$ euro mempunyai peluang $5\,\%$
        rusak total setiap tahun. Hitunglah premi asuransi yang adil,
        dan harapan laba penjamin per polis
        bila ia menarik $60$ euro.
\end{enumerate}

\medskip
\noindent\textbf{Bagian II --- Membagi taruhannya.}
Kedudukannya $5$--$3$, yang lebih dahulu mencapai $6$ mengambil $64$ pistole,
dan setiap ronde adalah lemparan koin yang setimbang.
\begin{enumerate}[resume]
  \item Dua jawaban keliru yang menggoda: bagilah $5 : 3$
        (sebanding dengan ronde yang dimenangkan), atau berikan semuanya kepada
        yang unggul. Katakan dalam satu kalimat untuk masing-masing apa yang
        tidak adil padanya.
  \item Metode Fermat: pemain A butuh $1$ kemenangan lagi, B butuh
        $3$; bayangkan bermain tepat $3$ ronde lagi apa pun yang terjadi.
        Daftarkan $8$ hasil yang sama-sama mungkin, tandai siapa yang menang
        rangkaiannya pada masing-masing, lalu simpulkan pembagian adil $64$
        pistole itu.
  \item Butir halus yang harus dibela Fermat: permainan yang sesungguhnya
        mungkin berhenti setelah satu ronde --- lalu mengapa bernalar atas
        semua $3$ ronde khayalan itu tetap sah?
  \item Metode Pascal, mundur dari ujungnya: hitunglah bagian adil A pada
        kedudukan (khayalan) $5$--$5$, lalu $5$--$4$, lalu $5$--$3$, setiap
        kali dengan merata-ratakan kedua masa depan yang sama-sama mungkin.
        Bandingkan dengan pertanyaan 7.
  \item Pertanyaan yang sama pada kedudukan $4$--$3$: berapa ronde khayalan,
        hasil mana yang memberi rangkaiannya kepada B, dan berapa pembagian
        adilnya?
  \item Nyatakan asas yang dibagi kedua metode itu --- nilai adil sebuah
        masa depan yang tak pasti adalah \emph{\omterm{def:g11:prob:expectation}{nilai harapannya}} ---
        lalu sebutkan dua industri modern yang seluruhnya berjalan atasnya
        (pertanyaan 5 memperlihatkan salah satunya).
  \item Sebuah pemeriksaan kewarasan untuk setiap kedudukan: jelaskan, dengan
        kelinearan \omterm{def:g11:prob:expectation}{nilai harapan}
        (\cref{prop:g11:prob:linearity}), mengapa bagian adil kedua pemain
        harus selalu berjumlah tepat sama dengan taruhan
        $64$ pistole itu.
\end{enumerate}

\medskip
\noindent\textbf{Bagian III --- Sulap kelinearan.}
\begin{enumerate}[resume]
  \item Temukan kembali $\E(\text{jumlah dua dadu}) = 7$ hanya dalam
        satu baris dengan kelinearan (\cref{prop:g11:prob:linearity}), tanpa
        tabel $36$ sel. Untuk \emph{\omterm{def:g11:prob:variance}{ragam}} jumlahnya,
        aturan $V(X + Y) = V(X) + V(Y)$ menuntut dadunya saling
        bebas (diterima tanpa bukti): hitunglah.
  \item Pesta penitipan topi: $n$ tamu melemparkan topinya ke dalam
        satu tumpukan lalu masing-masing mengambil satu kembali secara acak.
        Misalkan $X$ membilang tamu yang memperoleh topinya sendiri. Untuk satu
        tamu tertentu, berapa peluang memperoleh topinya sendiri? Dengan
        menjumlahkan $n$ peluang itu (kelinearan!), hitunglah $\E(X)$.
        Jawabannya semestinya mengejutkan: nilainya tidak bergantung pada $n$.
  \item Periksalah $\E(X) = 1$ secara paksa untuk $n = 2$ dan
        $n = 3$ (daftarkan $2$ dan $6$ penugasan topi yang sama-sama mungkin
        lalu rata-ratakan cacahnya).
  \item Perolehan topi para tamu jauh dari saling bebas (jika
        $n - 1$ tamu memperoleh topinya sendiri, maka tamu terakhir pun
        pasti memperolehnya). Di bagian mana persisnya perhitungan pertanyaan 14
        \emph{tidak} menuntut kebebasan? Nyatakan kekuatan super kelinearannya.
\end{enumerate}

\medskip
\noindent\textbf{Bagian IV --- Apa yang tak terlihat oleh \omterm{def:g11:prob:expectation}{nilai harapan}.}
\begin{enumerate}[resume]
  \item Dua kartu gosok seharga $5$ euro: kartu A memenangkan $10$ euro dengan
        peluang $\frac12$; kartu B memenangkan $1\,000$ euro dengan
        peluang $0.005$. Periksalah bahwa kedua keuntungannya mempunyai
        \omterm{def:g11:prob:expectation}{nilai harapan} yang sama, lalu hitunglah kedua \omterm{def:g11:prob:variance}{ragamnya}. Kartu
        mana yang sungguh dibeli orang, dan apa sebenarnya yang mereka beli?
  \item Permainan St Petersburg: lempar koin sampai angka pertama muncul;
        jika angka pertamanya muncul pada lemparan ke-$n$ kamu menerima
        $2^n$ euro. Hitunglah sumbangan setiap $n$ pada
        \omterm{def:g11:prob:expectation}{nilai harapannya} dan \omterm{def:g11:prob:expectation}{nilai harapannya} sendiri. Maukah kamu
        membayar $1\,000$ euro untuk bermain sekali? Apa yang diajarkan
        benturan itu tentang \omterm{def:g11:prob:expectation}{nilai harapan}?
  \item Asuransi dari sisi nasabah: membayar $20$ euro, atau
        menghadapi risiko $1\,\%$ kehilangan $1\,000$? Bandingkan kedua
        \omterm{def:g11:prob:expectation}{nilai harapannya}, lalu jelaskan mengapa membeli asuransi tetap dapat
        masuk akal (apa yang tidak dapat ``dirata-ratakan'' oleh satu keluarga?).
  \item Penutup --- riwayat hidup \omterm{def:g11:prob:expectation}{nilai harapan}, satu kalimat untuk setiap
        babnya: lahir pada tahun 1654 untuk membagi taruhan; kekuatan supernya,
        yaitu kelinearan, yang menjumlahkan \omterm{def:g11:prob:expectation}{nilai harapan} tanpa menuntut
        kebebasan (topi tadi); titik butanya, yaitu risiko, yang diukur oleh
        \omterm{def:g11:prob:variance}{ragam} (kartu gosok, St Petersburg, asuransi); dan penobatannya tahun
        berikutnya, ketika hukum bilangan besar menjelaskan tepatnya mengapa
        bandar selalu menang.

\end{enumerate}
\end{problem}
